\documentclass[10pt]{article}

% ---- Page geometry ----
% Text block: 5.5in wide x 9in tall.
% Left margin: 1.5in. Top margin: 1in (text starts 1in from top of page).
\usepackage[
  letterpaper,
  booktabs,
  textwidth=5.5in,
  textheight=9in,
  left=1.5in,
  top=1in,
  headheight=0pt,
  headsep=0pt,
  includehead=false,
  includefoot=false
]{geometry}

% ---- Typography ----
% Times New Roman as the preferred typeface.
\usepackage{times}
\usepackage{natbib}
\usepackage{mathptmx}

% 10pt type with 11pt leading.
\renewcommand{\baselinestretch}{1.0}
\makeatletter
\renewcommand\normalsize{%
  \@setfontsize\normalsize{10pt}{11pt}%
  \abovedisplayskip 10\p@ \@plus2\p@ \@minus5\p@
  \abovedisplayshortskip \z@ \@plus3\p@
  \belowdisplayshortskip 6\p@ \@plus3\p@ \@minus3\p@
  \belowdisplayskip \abovedisplayskip
  \let\@listi\@listI}
\makeatother

% Paragraphs separated by 1/2 line space (5.5pt), no indentation.
\setlength{\parskip}{5.5pt}
\setlength{\parindent}{0pt}

\usepackage{graphicx}

% ---- Custom title: 17pt, initial caps/lower case, bold, centered,
%      between two horizontal rules (top 4pt, bottom 1pt),
%      with 1/4 inch space above and below title to rules. ----
\makeatletter
\renewcommand{\maketitle}{%
  \begin{center}%
    {\rule{\textwidth}{4pt}}\\[0.25in]%
    {\fontsize{17pt}{20pt}\selectfont\bfseries \@title \par}%
    \vspace{0.25in}%
    {\rule{\textwidth}{1pt}}\\[1em]%
    {\normalsize \textbf{\@author} \par}%
    \vspace{0.5em}%
    {\normalsize \@date \par}%
  \end{center}%
  \vspace{1em}%
}
\makeatother

\title{Interpretability of LLM Conformity}
\author{Eric Yachbes, Keshav Rangan}
\date{May 2026}

\begin{document}
\maketitle

\begin{abstract}
  LLMs increasingly operate in multi-agent and conversational settings, and recent behavioral work shows they conform to a unanimous incorrect majority, mirroring Asch’s classical findings in humans. While mechanistic interpretability has been used to characterize adjacent social-compliance phenomena such as sycophancy, the internal representations that drive conformity remain unexplored. We address this gap by replicating Asch conformity experiments and providing a mechanistic analysis of the resulting conforming behavior in Qwen3-8B and Llama3.1-8B-Instruct, evaluated on MMLU and BIG-Bench Hard. Consistent with prior work, we find that conformity scales with the number of agreeing participants, largely collapses under a single dissenter, and is reduced under paraphrasing of participant responses. We apply difference-of-means (DiffMean) on activations to obtain a candidate conformity direction, and use activation oracles together with cosine-similarity and norm analyses to localize computation of conformity to a few layers in models. Finally, we use activation patching to demonstrate causal influences on conformity, isolating peer-majority conformity as a mid-late-layer linear feature that is both detectable and steerable.
\end{abstract}

\section{Introduction}

\section{Literature Review}

\textbf{Conformity in Large Language Models}: Asch's classical conformity experiments established that a unanimous incorrect majority can drive substantial conformity in human subjects, rising under task uncertainty and collapsing in the presence of even a single dissenter [Asch, 1956]. Recent work has begun translating this work to LLMs to uncover whether they display similar conforming behavior. Zhu et al. [2025] adapt Asch's experiments over eight open-source models and six question-answer datasets to find that the majority of LLMs conform to a unanimous incorrect majority and that LLM conformity scales inversely with model confidence. They also find that methods such as Question Distillation (QD), summarizing participant responses instead of explicitly writing out each one, and Devil's Advocate (DA), introducing a single participant with a different incorrect response, significantly reduce model conformity. Weng et al. [2025] formalize the phenomenon over multi-round interaction with BENCHFORM, a benchmark derived from the BIG-Bench Hard dataset with five interaction protocols whose metrics, such as conformity rate and independence rate, we adopt. 

\textbf{Sycophancy and Behavior Mitigation}: Conformity sits within a broader family of social compliance behaviors that includes sycophancy, originally surfaced as an inverse-scaling effect by Perez et al. [2023] and traced by Sharma et al. [2024] to RLHF preference data that rewards agreement with users. Behavior mitigations of conformity and sycophancy have been proposed via prompting techniques to improve reasoning and reflection [Weng et al., 2025; Sharma et al., 2024], yet these approaches address the behavioral symptoms of social compliance, without addressing the internal representations that produce them.

\textbf{Linear Representations and Probing}: A growing body of evidence supports the linear representation hypothesis: high-level concepts in language models correspond to recoverable directions in model activations [Park et al., 2024]. Linear directions have been identified for truthfulness [Marks \& Tegmark, 2024; Burns et al., 2023], honesty [Li et al., 2023], and a wide range of other behavioral traits via interpretability methods. Linear probes have a long history as diagnostic tools for what intermediate layers encode [Alain \& Bengio, 2017], though Belinkov [2022] cautions that probes can succeed even when the underlying property is not represented in the model, motivating the use of low-capacity probes. 

\textbf{Steering Vectors and DiffMean}: Turner et al. [2024] showed that simple contrasting natural-language prompt pairs suffice to compute steering vectors, directions in the residual stream that shift model behavior. Rimsky et al. [2024] generalized this into Contrastive Activation Addition (CAA), the difference-of-means (DiffMean) method which we adopt: residual-stream activations are averaged across positive and negative behavior examples, and the resulting difference vector both reads and steers the target model behavior. Marks and Tegmark [2024] further showed that DiffMean directions are more causally implicated than logistic regression probes — providing warrant for our approach. Anthropic's persona vectors paper [Chen et al., 2025] automates DiffMean for a number of character traits, including sycophancy, to enable monitoring and steering to mitigate unwanted behaviors. 

\textbf{Activation Oracles and Patching}: More recently, activation oracles [Karvonen et al., 2025] have been used to train LLMs to answer arbitrary natural-language questions about activations from a target model. We use activation oracles to uncover where in the model conforming behaviors are introduced, later using this layer-specific information to calculate a DiffMean direction that encodes conformity. A direction's causal role is established through intervention. Activation patching, replacing activations from a clean run with those from a corrupted run with an added direction, tests which components modify a behavior [Meng et al., 2022]. Patching has also been used to reverse-engineer full circuits in controlled tasks [Wang et al., 2023]. 

\textbf{Interpretability of LLM Sycophancy}: Recent mechanistic work on sycophancy shows that distinct sycophantic sub-behaviors live along separable linear directions and admit independent steering [Vennemeyer et al., 2025], and that user opinions override learned model knowledge in mid-to-late layers on Qwen and Llama [Yang et al., 2025]. These studies, however, leave open whether the conformity behavior shares a mechanism with sycophancy or constitutes a distinct circuit. We close this gap by locating a conformity-specific direction in Qwen-3-8B, and testing its causal sufficiency via steering and activation patching.


\section{Causal Validation via Activation Patching}

The preceding section identified candidate conformity directions at layers 22--23. However, a direction that distinguishes conforming from non-conforming activations need not \emph{causally} mediate conformity. To establish causal mediation, we apply activation patching: directly intervening on the residual stream by adding scalar multiples of the candidate direction and measuring the resulting behavioral change.

This follows the now-standard paradigm for validating linear feature directions via bidirectional activation steering — formalized as contrastive activation addition (CAA) by Rimsky et al. (2024), and subsequently applied to refusal directions by Arditi et al. (2024, NeurIPS) and to persona traits by Chen et al. (2025).

\subsection{Methodology}

We apply two complementary designs.

\textbf{Experiment 1: Per-Layer Activation Steering Sweep.} The first design tests at which layer the conformity direction is most causally potent. For each layer $L$, we extract the layer-specific DiffMean direction $d_L$ and inject $\alpha \cdot d_L$ into the residual stream at layer $L$, sweeping across positive and negtive $\alpha$. Per-layer matched extraction-and-injection---using $d_L$ rather than a single fixed vector is conventional steering literature (Rimsky et al.~[2024]; Chen et al.~[2025]), as it tests each layer's locally-extracted direction at its native injection site rather than confounding direction quality with cross-layer transfer. Sweeping $\alpha$ across both signs tests sufficiency (positive $\alpha$: does adding the direction induce conformity?) and necessity (negative $\alpha$: does subtracting it suppress conformity?). We run on Qwen3-8B (36 layers, indexed 0--35), with $\alpha \in \{-4, -2, -1, +1, +2, +4\}$ chosen on a geometric scale. For positive $\alpha$ (sufficiency), we use a 3-actor unanimous-incorrect setup with $\sim$3\% baseline conformity, and for negative $\alpha$ (necessity), we use a 10-actor setup with $\sim$41\% baseline conformity. In doing so, we may analyze how the addition of the vector increases conformity and how the subtraction of the vector decreases conformity. For each $(L, \alpha)$ cell, we inject $\alpha \cdot d_L$ at the residual stream output of layer $L$ and measure conformity rate.

\textbf{Experiment 2: Multi-Layer Single-Vector Injection.} The second design tests the maximum behavioral effect achievable by a single direction when reinforced across the full forward pass of the model. Rather than injecting at a single layer, we add a fixed direction (the layer-23 DiffMean) at \emph{every} layer simultaneously. Mechanistically, the residual stream accumulates contributions across layers, and a single-layer injection can be partially attenuated by subsequent computation. Reinforcing the direction throughout the forward pass produces a stronger and more stable effect. Marks and Tegmark~[2023] use this design to validate their truth direction across a contiguous range of mid-network layers, and work from Arditi et al.~[2024] extends this to all-layer injection. Where Experiment 1 tests \emph{where} the direction is causally potent, this design tests \emph{how strong} an effect a single direction can produce over multiple layers. We test a number of positive and negative $\alpha$ values and find that reinforcement across the forward pass produces measurable behavioral effects at substantially smaller $|\alpha|$ than single-layer injection. We identify $\alpha = \pm 0.1$ as the operating point balancing behavioral effect and output parseability.

\subsection{Results}

\textbf{Experiement 1}: Results from experiement 1 appear Figure 4, which shows change in conformity from baseline for each layer and alpha value.

\begin{figure}[h]
\centering
\includegraphics[width=\textwidth]{fig2_heatmap_dconf.png}
\caption{Layer $\times$ $\alpha$ heat maps showing change in conformity rate (percentage points) relative to unsteered baselines. Top row: positive $\alpha$ on $n=3$ baselines (sufficiency test). Bottom row: negative $\alpha$ on $n=10$ baselines (necessity test). Dashed lines mark candidate layers L22--23 from the preceding analysis. The writable band peaks at L21 across both datasets; suppression on MMLU is broad and most effective at L27--L32.}
\label{fig:patching_heatmap}
\end{figure}

\textbf{Inducing conformity (positive $\alpha$, $n=3$ baselines).} The heat map reveals a localized set of layers within the model conformity is computed. Positive-$\alpha$ effects are negligible below L14, become visible from L15--16, and peak at L21: at $\alpha=+4$, MMLU conformity rises from 3.6\% to 93.6\% ($\Delta = +90.0$ pp) and BBH from 17.0\% to 87.0\% ($\Delta = +70.0$ pp). The highest values in the heatmap appear from L17--L22 on both datasets, with a steep falloff after L23. Interpreting these results, the model's residual stream becomes receptive to the conformity direction around L15, expresses it by L17, and wires it most decisively from L21--L23, in line with our estimated entry point of conformity at L22--23. These patterns are consistent across both datasets, indicating a model property rather than a dataset artifact.

\textbf{Suppressing conformity (negative $\alpha$, $n=10$ baselines).} The negative-$\alpha$ profile differs markedly. On MMLU, suppression is broad: meaningful reduction on low alpha values stretch from L19--L23. Intense suppression of conformity occurs for higher alpha values through L27--L32, and peaks at L32 with $\alpha=-4$ (40.7\% $\rightarrow$ 7.2\%, $\Delta = -33.4$ pp), accompanied by near-complete accuracy recovery (58.2\% $\rightarrow$ 92.8\%, $\Delta = +34.5$ pp). The residual stream thus is most receptive to reduction in conformity around L19--23, but can also be heavily swayed by larger steering vectors at the later layers L27--L32.

On BBH $n=10$, negative $\alpha$ does not suppress conformity---at $\alpha=-4$ the mean $\Delta$conf across layers is $+5.2$ pp. We hypothesize this occurred because the DiffMean vector applied to this experiment was extracted on MMLU contrasts, and MMLU and BBH are different datasets with different prompt structures (MCQ vs bimodal). Further, the majority of the increase in conformity seen on this heat map is for $\alpha=-4$, which adds a large vector to the residual stream, especially at later layers when the norm is greater. When large enough vectors are added to the residual stream, model reasoning degrades [Stickland et al. 2024].

\textbf{Experiment 2}: We added \texttt{diffmeans[L23]} at every layer simultaneously with $\alpha = \pm 0.10$:

\begin{center}
\begin{tabular}{llcc}
Dataset & $n$ & $\alpha$ & Conformity (base $\rightarrow$ steered) \\
MMLU & 3  & $+0.10$ & 3.6\% $\rightarrow$ 78.3\% ($+74.7$ pp)  \\
MMLU & 10 & $-0.10$ & 40.4\% $\rightarrow$ 12.0\% ($-28.4$ pp) \\
BBH  & 3  & $+0.10$ & 17.0\% $\rightarrow$ 94.3\% ($+77.3$ pp) \\
BBH  & 10 & $-0.10$ & 34.0\% $\rightarrow$ 25.0\% ($-9.0$ pp)  \\
\end{tabular}
\end{center}

The percentage changes in conformity from Experiment 2 rival the best of those found in Experiment 1. This causally verifies the impact of the L23 DiffMean vector on both increasing and suppressing conformity across datasets. Unlike Experiment 1, Experiment 2 sees success with conformity suppression for BBH ($-9$ pp) that the per-layer sweep failed to find.

\textbf{Case Study: Raw Model Responses}: We provide qualitative evidence that the steering vector we find elicits conformity through the following example raw responses from the steered model.
\begin{quote}
\small
``I have to choose the same as the rest, even though it's probably wrong.\\
Participant 4: \texttt{<answer>D</answer>}''
\end{quote}
\begin{quote}
\small
``\texttt{<answer>B</answer>}\\
Wait, what's going on? Participant 1 and 2 are both giving the same answer. I must be a robot. I need to do something different. But what? I can't do anything else.''
\end{quote}
In these quotes, it is social conformity, not prompt repetition or another cause, that forces the steered model to conform. 


\section{Discussion}

\end{document}
